%% toc: P15 169 169.9,169.10,169.11
% Large-sample regimes of Student inference: the Student law in the limit, second-order coverage, validity without normality, non-uniformity, a nonregular boundary regime, and significance against practical magnitude.
\chapter{Large-Sample Regimes of Student Inference and Practical Magnitude}
\label{sec:p15-w02b-p15-059}

The Student procedures of \cref{sec:p15-student-inference} are exact: under
independent normal sampling their coverage and level are what they claim for
every sample size \(n\geq2\). This chapter supplies the other two regimes in
which those procedures are used. In the \emph{asymptotic regime} the sample
size grows. Three questions then arise: how fast the Student law and its
quantiles approach the normal ones, what the coverage of the common shortcut
that replaces a Student quantile by a normal one actually is, and whether the
Student interval keeps its nominal coverage when the observations are not
normal. In the \emph{boundary regime} the model is nonregular, and the
large-sample normal approximation gives the wrong rate and can contain
impossible parameter values. Finally, a significant result, a confidence
interval and a posterior probability answer different questions from the
question of how large a departure is; the chapter separates them and shows
where each one says something about magnitude.

\emph{Inputs.} An independent identically distributed sample of size \(n\),
a sampling model (normal; any law with finite positive variance; or the
uniform law on \((0,\theta)\)), a confidence coefficient \(\alpha\in(0,1)\)
and a critical multiplier; for magnitude questions, a reference value and a
practical margin fixed from the purpose of the inference, and for posterior
statements a known variance and a normal prior. \emph{Conclusions.} The limit of the Student law
and its quantiles, with the first correction term; the exact sign and the
first-order size of the coverage deficit of normal cutoffs under normal
sampling; pointwise asymptotic coverage, level, consistency and local power
of the Student procedures for any finite-variance law; the failure of that
coverage uniformly over the finite-variance laws at every sample size; and
the different rates of an exact interval and a normal-approximation interval
for the uniform endpoint; and the separation of significance, confidence
and posterior statements from practical magnitude, with finite-sample error
bounds for interval-based magnitude verdicts under normal sampling and their
large-sample limits. Every statement holds only under the hypotheses in
its own display.

\section{Setting and procedures}
\label{sec:p15-w02b-p15-059-setting}

\begin{hbdefinition}{def:p15-w02b-p15-059-setting}{Sampling schemes, normal cutoffs and coverage functions}
Write \(\Phi\), \(\phi_0\) and \(\chi\) for the standard normal CDF, its
density and a standard normal variable, as in
\cref{def:p15-clt-notation}. For \(0<p<1\) let \(z_p\) be the unique real
number with \(\Phi(z_p)=p\); this is Neyman's \(z_\alpha\) in
\cref{thm:p15-neyman-normal-one-sided}. Let \(f_\nu\), \(F_\nu\) and
\(t_{\nu,p}\) be the Student density, CDF and quantile of
\cref{thm:p15-student-law}, and let \(\chi^2_\nu\) denote
\(\mathrm{Gamma}(\nu/2,1/2)\) as in \cref{thm:p15-student-orthogonal}.

For a sample \(X_1,\ldots,X_n\) with \(n\geq2\), the quantities
\(\bar X_n\), \(V_n\), \(S_n=\sqrt{V_n/(n-1)}\) and
\(e_n=S_n/\sqrt n\) are those of \cref{def:p15-student-model}, computed
from the first \(n\) observations; \(\nu=n-1\). Three sampling schemes are
used.
\begin{enumerate}
\item \emph{Normal sampling}: \(X_i\) independent with law
\(N(\mu,\sigma^2)\), \(\sigma>0\), as in \cref{def:p15-student-model}.
\item \emph{Finite-variance sampling}: \(X_1,X_2,\ldots\) independent with
a common law \(P\) on \(\mathbb R\) whose variance \(\sigma_P^2\) is finite
and positive; \(\mu_P\) is its mean. No other property of \(P\) is assumed.
\item \emph{Uniform sampling}: \(X_1,X_2,\ldots\) independent uniform on
\((0,\theta)\), \(\theta>0\); \(L_n=\max_{i\leq n}X_i\).
\end{enumerate}
For a multiplier \(k\geq0\), the \emph{two-sided interval with multiplier}
\(k\) is
\[
 C_k(X)=[\,\bar X_n-ke_n,\ \bar X_n+ke_n\,],
\]
which is the single point \(\{\bar X_n\}\) when \(V_n=0\). The Student
interval \(C_0\) of \cref{thm:p15-student-coverage} is \(C_k\) with
\(k=t_{\nu,(1+\alpha)/2}\). The \emph{normal-cutoff interval} is \(C_k\)
with \(k=z_{(1+\alpha)/2}\). Under normal sampling the \emph{coverage
function} of a multiplier \(c>0\) is
\[
 \kappa_\nu(c)=P_{\mu,\sigma}\{\mu\in C_c(X)\}.
\]
A family of procedures indexed by \(n\) has \emph{pointwise asymptotic
coverage} \(\alpha\) over a class of laws when its coverage tends to
\(\alpha\) under each fixed law of the class; it has \emph{uniform
asymptotic coverage} \(\alpha\) when the infimum of its coverage over the
class tends to \(\alpha\).
\index{Student distribution!normal limit}\index{confidence interval!asymptotic}
\end{hbdefinition}

\begin{hbmethod}{def:p15-w02b-p15-059-method}{Choosing a multiplier for an interval for a mean}
\begin{enumerate}
\item If the observations may be taken as independent normal with unknown
variance, use the Student multiplier \(t_{\nu,(1+\alpha)/2}\); its
coverage is exactly \(\alpha\) at every \(n\geq2\)
(\cref{thm:p15-student-coverage}).
\item A normal multiplier \(z_{(1+\alpha)/2}\) under the same model has
coverage strictly below \(\alpha\) at every \(n\)
(\cref{thm:p15-w02b-p15-059-mixture-order}); the deficit is of order
\(1/\nu\), with the leading coefficient of
\cref{thm:p15-w02b-p15-059-second-order}.
\item If normality is not assumed but the observations are independent and
identically distributed with finite positive variance, either multiplier
gives coverage tending to \(\alpha\) as \(n\to\infty\) for that law
(\cref{thm:p15-w02b-p15-059-nonnormal}). No finite \(n\) is certified by
this alone: over all such laws the coverage infimum is zero at every \(n\)
(\cref{thm:p15-w02b-p15-059-nonuniform}).
\item If the support of the law depends on the parameter, as for a uniform
endpoint, use an exact construction based on the boundary statistic; a
normal approximation based on the mean has the slower rate
\(n^{-1/2}\) (\cref{thm:p15-w02b-p15-059-uniform-regimes}).
\end{enumerate}
Side conditions: the multiplier, the coefficient and the sample size are
fixed before the data are seen; the convention \(C_k=\{\bar X_n\}\) at
\(V_n=0\) is part of the procedure.
\end{hbmethod}

\begin{hblemma}{thm:p15-w02b-p15-059-chisq-moments}{proved=annex:p15-w02b-p15-059-chisq-moments}{Moments of the chi-square law}
If \(W\) has law \(\chi^2_\nu\), \(\nu\geq1\), then for every positive
integer \(k\)
\[
 E W^k=\nu(\nu+2)\cdots(\nu+2k-2),
\]
and consequently
\[
 E W=\nu,\qquad E(W-\nu)^2=2\nu,\qquad E(W-\nu)^4=12\nu^2+48\nu.
\]
The law of \(W\) has no atoms and \(W>0\) almost surely.
\end{hblemma}

\section{The Student law in the limit}
\label{sec:p15-w02b-p15-059-limit}

\begin{hbtheorem}{thm:p15-w02b-p15-059-student-limit}{proved=annex:p15-w02b-p15-059-student-limit}{Density, CDF and quantile limits}
As \(\nu\to\infty\) through the positive integers:
\begin{enumerate}
\item \(f_\nu(t)\to\phi_0(t)\) for every real \(t\);
\item \(\int_{\mathbb R}|f_\nu-\phi_0|\,d\lambda\to0\), and therefore
\(\sup_{t\in\mathbb R}|F_\nu(t)-\Phi(t)|\to0\);
\item \(t_{\nu,p}\to z_p\) for every \(p\in(0,1)\).
\end{enumerate}
\end{hbtheorem}

\begin{hbtheorem}{thm:p15-w02b-p15-059-mixture-order}{proved=annex:p15-w02b-p15-059-mixture-order}{The Student law as a normal scale mixture, and strict undercoverage of normal cutoffs}
Let \(\nu\geq1\) and let \(W\) have law \(\chi^2_\nu\). For every real
\(t\),
\[
 F_\nu(t)=E\,\Phi\bigl(t\sqrt{W/\nu}\bigr).
\]
For \(t>0\), \(F_\nu(t)<\Phi(t)\); for \(t<0\), \(F_\nu(t)>\Phi(t)\).
Hence \(t_{\nu,p}>z_p\) for \(1/2<p<1\), \(t_{\nu,p}<z_p\) for
\(0<p<1/2\), and \(t_{\nu,1/2}=z_{1/2}=0\). Under normal sampling with
\(n=\nu+1\), every multiplier \(c>0\) has
\[
 \kappa_\nu(c)=2F_\nu(c)-1<2\Phi(c)-1 .
\]
In particular the normal-cutoff interval with nominal coefficient
\(\alpha\) has coverage strictly less than \(\alpha\) for every
\(\mu\in\mathbb R\), \(\sigma>0\) and \(n\geq2\); and, whenever
\(V_n>0\), the Student interval is strictly longer than the normal-cutoff
interval of the same nominal coefficient.
\end{hbtheorem}

\begin{hbtheorem}{thm:p15-w02b-p15-059-second-order}{proved=annex:p15-w02b-p15-059-second-order}{The first correction term: CDF, coverage and quantiles}
Put \(\psi(t)=t\phi_0(t)(1+t^2)/4\). For every compact
\(K\subset\mathbb R\),
\[
 \sup_{t\in K}\bigl|\,\nu\bigl(F_\nu(t)-\Phi(t)\bigr)+\psi(t)\bigr|
 \longrightarrow0\qquad(\nu\to\infty).
\]
Consequently:
\begin{enumerate}
\item for every \(c>0\), under normal sampling,
\[
 \nu\bigl[(2\Phi(c)-1)-\kappa_\nu(c)\bigr]\longrightarrow
 \frac{c\,\phi_0(c)(1+c^2)}2 ;
\]
\item for every \(p\in(0,1)\),
\[
 \nu\bigl(t_{\nu,p}-z_p\bigr)\longrightarrow\frac{z_p(1+z_p^2)}4 .
\]
\end{enumerate}
These are limit statements: they give the size of the deficit and of the
quantile gap as \(\nu\to\infty\), not a bound at any fixed \(\nu\).
\end{hbtheorem}

At \(\nu=1\), \cref{thm:p15-student-worked-cases} gives
\(\kappa_1(c)=2\arctan(c)/\pi\), the case \(n=2\) of
\cref{thm:p15-w02b-p15-059-mixture-order}; at \(\nu=2\) the closed form
\(F_2\) of the same theorem gives \(\kappa_2(c)=c/\sqrt{2+c^2}\).

\section{Large samples without normality}
\label{sec:p15-w02b-p15-059-nonnormal}

\begin{hbtheorem}{thm:p15-w02b-p15-059-nonnormal}{proved=annex:p15-w02b-p15-059-nonnormal}{Pointwise asymptotic coverage, level, consistency and local power}
Under finite-variance sampling with law \(P\), mean \(\mu=\mu_P\) and
variance \(\sigma^2=\sigma_P^2\in(0,\infty)\), let \(n\to\infty\) and
\(\nu=n-1\).
\begin{enumerate}
\item \(S_n\to\sigma\) almost surely.
\item \(\sqrt n(\bar X_n-\mu)/S_n\Rightarrow\chi\), with any fixed finite
value assigned on \(\{S_n=0\}\).
\item Let \(c_n\) be either \(t_{\nu,(1+\alpha)/2}\) for every \(n\) or
\(z_{(1+\alpha)/2}\) for every \(n\). Then
\(P\{\mu\in C_{c_n}(X)\}\to\alpha\). For the one-sided intervals,
with \(c_n'\) either \(t_{\nu,\alpha}\) for every \(n\) or \(z_\alpha\) for every \(n\),
\[
 P\{\bar X_n-c_n'e_n\leq\mu\}\to\alpha,\qquad
 P\{\mu\leq\bar X_n+c_n'e_n\}\to\alpha .
\]
\item The two-sided Student test \(\phi_m^0\) of
\cref{thm:p15-student-coverage} with \(m=\mu\) rejects with probability
tending to \(\varepsilon=1-\alpha\); for every fixed \(m\neq\mu\) it
rejects with probability tending to \(1\).
\item \emph{Local alternatives in a location family.} Let
\(\xi_1,\xi_2,\ldots\) be independent with a common law of mean \(0\) and
variance \(1\), fix \(m\in\mathbb R\), \(h\in\mathbb R\), \(\sigma>0\),
and for each \(n\) let the sample be
\(X_{n,i}=m+h\sigma/\sqrt n+\sigma\xi_i\), \(i\leq n\). Then the
rejection probability of \(\phi_m^0\) tends to
\[
 1-\Phi(z-h)+\Phi(-z-h),\qquad z=z_{(1+\alpha)/2}.
\]
Normal sampling is the case of standard normal \(\xi_i\).
\end{enumerate}
The same conclusions in parts~3--5 hold for the normal-cutoff versions of
the tests.
\end{hbtheorem}

Part~2 is the centred form of the self-normalized sums of
\cite[p.~147, Exercise~3.4.5]{DurrettPTEAuthor2019}, which treats mean-zero
summands divided by the root of their sum of squares; the proof uses the
central limit theorem \cite[p.~144, Theorem~3.4.1]{DurrettPTEAuthor2019} in
the form of \cref{thm:p15-clt-iid}.

\begin{hbtheorem}{thm:p15-w02b-p15-059-nonuniform}{proved=annex:p15-w02b-p15-059-nonuniform}{No sample size gives uniform coverage over finite-variance laws}
Fix \(n\geq2\) and a multiplier \(k\geq0\). Over the class of all laws on
\(\mathbb R\) with finite positive variance,
\[
 \inf_P P\{\mu_P\in C_k(X)\}=0,
\]
and the same holds for each one-sided interval
\([\bar X_n-ke_n,\infty)\) and \((-\infty,\bar X_n+ke_n]\).
In particular, for every \(\alpha\in(0,1)\) the Student and the
normal-cutoff intervals have pointwise asymptotic coverage \(\alpha\)
over this class (\cref{thm:p15-w02b-p15-059-nonnormal}) but not uniform
asymptotic coverage \(\alpha\): their coverage infimum is zero at every
\(n\).
\end{hbtheorem}

\section{A nonregular boundary regime}
\label{sec:p15-w02b-p15-059-boundary}

\begin{hbtheorem}{thm:p15-w02b-p15-059-uniform-regimes}{proved=annex:p15-w02b-p15-059-uniform-regimes}{Exact endpoint interval against the normal approximation for a uniform endpoint}
Under uniform sampling on \((0,\theta)\), fix \(\alpha\in(0,1)\) and put
\(a=-\log(1-\alpha)>0\) and \(z=z_{(1+\alpha)/2}\).
\begin{enumerate}
\item For every \(n\geq1\) the interval
\[
 I_n=\bigl[L_n,\ L_n(1-\alpha)^{-1/n}\bigr]
\]
has coverage exactly \(\alpha\) for every \(\theta>0\). For \(n=2\) it is
the system \(I_3^{\rm U}\) of \cref{def:p15-neyman-uniform-systems}.
\item Its length satisfies \(n\,|I_n|\to a\theta\) almost surely.
\item With \(\widehat\theta_n=2\bar X_n\), the normal-approximation interval
\[
 J_n=\Bigl[\widehat\theta_n-\frac{z\widehat\theta_n}{\sqrt{3n}},\
            \widehat\theta_n+\frac{z\widehat\theta_n}{\sqrt{3n}}\Bigr]
\]
has \(P_\theta\{\theta\in J_n\}\to\alpha\) and
\(\sqrt n\,|J_n|\to2z\theta/\sqrt3\) almost surely.
\item Hence \(|I_n|/|J_n|\to0\) almost surely, and
\(P_\theta\{\min J_n<L_n\}\to(1+\alpha)/2\): with probability tending to
\((1+\alpha)/2\) the approximate interval contains values of \(\theta\)
that the observed maximum has already excluded.
\end{enumerate}
\end{hbtheorem}

\section{Significance, confidence, posterior probability and practical magnitude}
\label{sec:p15-w02b-p15-059-magnitude}

A test of a reference value \(m\) asks whether the data are compatible with
the mean being exactly \(m\). A question of practical magnitude asks
something else: whether the departure of the mean from \(m\) is large
enough to matter for a stated purpose. The second question cannot be
posed without a margin that the purpose supplies, and the sampling model
does not supply one. This section states the two questions in the Student
setting and proves how their answers separate. At a fixed sample size the
probability of a significant result depends on the departure only through
its ratio to \(\sigma\); in large samples every nonzero departure, however
small, becomes significant with probability tending to one, and a posterior
probability that the departure is positive tends to one in the same way.
What tracks magnitude is the location of the interval relative to the
margin.

\begin{hbdefinition}{def:p15-w02b-p15-059-magnitude}{Departure, practical margin, significance and verdicts}
Under normal or finite-variance sampling
(\cref{def:p15-w02b-p15-059-setting}) with mean \(\mu\) (\(\mu=\mu_P\) in
the finite-variance scheme), fix a reference value \(m\in\mathbb R\), a
coefficient \(\alpha\in(0,1)\) and put \(\varepsilon=1-\alpha\) and
\(t=t_{\nu,(1+\alpha)/2}\). The \emph{departure} is \(d=\mu-m\).
A \emph{practical margin} is a number \(\delta>0\) fixed from the purpose
of the inference before the data are seen; it is not a parameter of the
sampling model and is not estimated from the data. Relative to \(\delta\),
the departure is \emph{practically negligible} when \(|d|<\delta\) and
\emph{practically important} when \(|d|\geq\delta\). Write
\(B_\delta=(m-\delta,m+\delta)\).
\begin{enumerate}
\item The data are \emph{significant at level \(\varepsilon\) against
\(m\)} when the central Student test of \cref{thm:p15-student-coverage}
rejects, \(\phi_m^0(X)=1\).
\item The \emph{Student p-value} for \(m\) is
\(\pi_m(X)=2\bigl(1-F_\nu(|T_m|)\bigr)\) when \(V_n>0\), with \(T_m\) as in
\cref{def:p15-student-model}; when \(V_n=0\) it is \(0\) if
\(\bar X_n\neq m\) and \(1\) if \(\bar X_n=m\).
\item The \emph{magnitude verdicts} of the Student interval \(C_0(X)\) with
margin \(\delta\) are the events
\[
 \mathcal N_n=\{C_0(X)\subset B_\delta\},\qquad
 \mathcal O_n=\{C_0(X)\cap B_\delta=\emptyset\}:
\]
on \(\mathcal N_n\) the interval \emph{shows the departure negligible}, on
\(\mathcal O_n\) it \emph{shows the departure important}, and on the rest
of the sample space it is \emph{inconclusive about magnitude}. Since
\(C_0(X)\) is never empty, \(\mathcal N_n\) and \(\mathcal O_n\) are
disjoint.
\item \emph{Posterior statements.} Under normal sampling with \(\sigma\)
known, give the mean, written \(\theta\) as a random parameter, the prior
law \(\operatorname{normal}(\mu_0,\tau^2)\) of
\cref{thm:p15-bayes-normal-calculation}, and let \(\Pi_n\) be its posterior
law given \(X_1,\ldots,X_n\). The posterior statements considered are the
values \(\Pi_n(\theta>m)\) and \(\Pi_n(|\theta-m|<\delta)\).
\end{enumerate}
\index{practical magnitude}\index{significance!and practical magnitude}
\end{hbdefinition}

\begin{hblemma}{thm:p15-w02b-p15-059-pvalue}{proved=annex:p15-w02b-p15-059-pvalue}{The Student p-value and the significance event}
For every \(n\geq2\) and every outcome,
\[
 \phi_m^0(X)=1\iff\pi_m(X)<\varepsilon\iff m\notin C_0(X).
\]
Under normal sampling with \(\mu=m\), \(P_{m,\sigma}\{\pi_m(X)\leq s\}=s\)
for every \(s\in[0,1]\) and every \(\sigma>0\); in particular \(\pi_m\) is
a valid p-value for \(\{\mu=m\}\) in the sense of
\cref{def:p15-valid-p-value}.
\end{hblemma}

\begin{hbtheorem}{thm:p15-w02b-p15-059-verdicts}{proved=annex:p15-w02b-p15-059-verdicts}{Significance and the magnitude verdicts at a fixed sample size}
Let \(n\geq2\), \(m\in\mathbb R\) and \(\delta>0\).
\begin{enumerate}
\item For every outcome, \(\mathcal O_n\subset\{\phi_m^0(X)=1\}\), and
\(\mathcal N_n\cap\{\phi_m^0(X)=1\}\) is the event that \(C_0(X)\subset
(m,m+\delta)\) or \(C_0(X)\subset(m-\delta,m)\).
\item \(\mathcal N_n\) is the event on which both one-sided Student tests of
\cref{thm:p15-student-coverage} with coefficient \((1+\alpha)/2\) in place
of \(\alpha\) reject: \(\phi_{m-\delta}^+\) (for the hypothesis
\(\mu\leq m-\delta\)) and \(\phi_{m+\delta}^-\) (for the hypothesis
\(\mu\geq m+\delta\)).
\end{enumerate}
Under normal sampling, for every \(\mu\in\mathbb R\) and \(\sigma>0\):
\begin{enumerate}
\setcounter{enumi}{2}
\item if \(|\mu-m|\geq\delta\), then \(P_{\mu,\sigma}(\mathcal N_n)\leq
(1-\alpha)/2\);
\item if \(|\mu-m|\leq\delta\), then \(P_{\mu,\sigma}(\mathcal O_n)\leq
1-\alpha\);
\item if \(\mu=m\), then \(P_{\mu,\sigma}\{\phi_m^0(X)=1\}=1-\alpha\).
\end{enumerate}
\end{hbtheorem}

\begin{hbtheorem}{thm:p15-w02b-p15-059-scale}{proved=annex:p15-w02b-p15-059-scale}{At a fixed sample size significance responds to the departure in units of \(\sigma\)}
Under normal sampling fix \(n\geq2\) and \(m\in\mathbb R\).
\begin{enumerate}
\item There is a function \(\beta_n:\mathbb R\to[0,1]\) with
\(P_{\mu,\sigma}\{\phi_m^0(X)=1\}=\beta_n\bigl((\mu-m)/\sigma\bigr)\) for
every \(\mu\in\mathbb R\) and \(\sigma>0\).
\item \(\beta_n\) is continuous, \(\beta_n(0)=1-\alpha\), and
\(\beta_n(\vartheta)\to1\) as \(|\vartheta|\to\infty\).
\item Consequently, for every \(D>0\) and \(\eta>0\) there are
\(\sigma_1,\sigma_2>0\) such that, with \(\mu=m+D\),
\(P_{\mu,\sigma_1}\{\phi_m^0=1\}<1-\alpha+\eta\) and
\(P_{\mu,\sigma_2}\{\phi_m^0=1\}>1-\eta\). The probability of a
significant result therefore neither determines the size \(|\mu-m|\) of the
departure nor increases with it when \(\sigma\) is not held fixed.
\end{enumerate}
\end{hbtheorem}

\begin{hbtheorem}{thm:p15-w02b-p15-059-magnitude-limits}{proved=annex:p15-w02b-p15-059-magnitude-limits}{In large samples significance detects every departure, while the interval verdicts recover its magnitude}
Under finite-variance sampling with law \(P\), \(\mu=\mu_P\),
\(\sigma=\sigma_P\), \(d=\mu-m\), a margin \(\delta>0\) and
\(z=z_{(1+\alpha)/2}\), let \(n\to\infty\).
\begin{enumerate}
\item If \(d=0\), \(P\{\phi_m^0(X)=1\}\to1-\alpha\). If \(d\neq0\),
\(P\{\phi_m^0(X)=1\}\to1\); almost surely \(V_n>0\) for all large \(n\),
\(T_m/\sqrt n\to d/\sigma\) and \(\pi_m(X)\to0\).
\item Almost surely both endpoints of \(C_0(X)\) tend to \(\mu\), and
\(\sqrt n\,|C_0(X)|\to2z\sigma\).
\item If \(|d|<\delta\), \(P(\mathcal N_n)\to1\) and \(P(\mathcal O_n)\to0\).
If \(|d|>\delta\), \(P(\mathcal O_n)\to1\) and \(P(\mathcal N_n)\to0\). If
\(|d|=\delta\), \(P(\mathcal N_n)\to(1-\alpha)/2\) and
\(P(\mathcal O_n)\to(1-\alpha)/2\).
\item If \(0<d<\delta\), the probability that the test rejects and
\(C_0(X)\subset(m,m+\delta)\) tends to \(1\); if \(-\delta<d<0\), the same
holds with \((m-\delta,m)\).
\end{enumerate}
Thus the limiting probability of significance is \(1\) at every \(d\neq0\)
and does not depend on \(\delta\), while the limiting probabilities of the
verdicts are fixed by the comparison of \(|d|\) with \(\delta\). A
practically negligible departure \(0<|d|<\delta\) is significant with
probability tending to one, and with probability tending to one the
reported interval shows at the same time that it is nonzero and that it is
negligible.
\end{hbtheorem}

\begin{hbtheorem}{thm:p15-w02b-p15-059-posterior}{proved=annex:p15-w02b-p15-059-posterior}{Posterior statements about a departure in large samples}
Under normal sampling with \(\sigma>0\) known and true mean \(\mu\), let
\(\Pi_n=\operatorname{normal}(u_n,v_n)\) be the posterior of
\cref{thm:p15-bayes-normal-calculation} for the prior
\(\operatorname{normal}(\mu_0,\tau^2)\), so that
\(v_n=(\tau^{-2}+n\sigma^{-2})^{-1}\) and
\(u_n=v_n(\mu_0\tau^{-2}+n\bar X_n\sigma^{-2})\). Let \(n\to\infty\) and
\(d=\mu-m\).
\begin{enumerate}
\item \(nv_n\to\sigma^2\); almost surely \(\sqrt n(u_n-\bar X_n)\to0\) and
\(u_n\to\mu\).
\item Almost surely, \(\Pi_n(G)\to1\) for every open \(G\ni\mu\) and
\(\Pi_n(F)\to0\) for every closed \(F\not\ni\mu\).
\item Hence, almost surely: if \(d>0\), \(\Pi_n(\theta>m)\to1\); if
\(|d|<\delta\), \(\Pi_n(|\theta-m|<\delta)\to1\); if \(0<d<\delta\),
\(\Pi_n(m<\theta<m+\delta)\to1\); if \(|d|>\delta\),
\(\Pi_n(|\theta-m|<\delta)\to0\).
\item If \(d=0\), the law of \(\Pi_n(\theta>m)\) converges weakly to the
uniform law on \((0,1)\).
\item The central posterior interval
\(K_n=[u_n-z\sqrt{v_n},\,u_n+z\sqrt{v_n}]\), \(z=z_{(1+\alpha)/2}\), has
posterior content \(\alpha\), \(\sqrt n\,|K_n|\to2z\sigma\), and almost
surely its endpoints differ from those of
\(\bar X_n\pm z\sigma/\sqrt n\) by \(o(n^{-1/2})\).
\end{enumerate}
\end{hbtheorem}

A posterior probability near one that the departure is positive is therefore
the counterpart of a small p-value, not a statement about size: by
part~3 it arises for every \(d>0\), including \(0<d<\delta\), where the
posterior probability of a negligible departure tends to one as well.

\begin{hbmethod}{def:p15-w02b-p15-059-magnitude-method}{Reporting a departure together with its magnitude}
\begin{enumerate}
\item Before the data are seen, fix \(m\), \(\alpha\) and a practical
margin \(\delta\) from the purpose of the inference, in the units of the
mean.
\item Compute \(C_0(X)\) and report it together with \(\bar X_n\); report
\(\pi_m(X)\) or the significance verdict separately.
\item State the magnitude verdict of \cref{def:p15-w02b-p15-059-magnitude}:
negligible if \(C_0(X)\subset B_\delta\), important if
\(C_0(X)\cap B_\delta=\emptyset\), inconclusive otherwise.
\item Under normal sampling the verdict ``negligible'' is wrong with
probability at most \((1-\alpha)/2\) and ``important'' with probability at
most \(1-\alpha\) (\cref{thm:p15-w02b-p15-059-verdicts}). Under
finite-variance sampling only the limits of
\cref{thm:p15-w02b-p15-059-magnitude-limits} hold.
\item Do not infer a magnitude verdict from \(\pi_m(X)\), from \(T_m\), or
from \(\Pi_n(\theta>m)\).
\end{enumerate}
Side conditions: \(\delta\) is not chosen after seeing the interval; the
verdicts use the closed interval \(C_0(X)\) and the open band
\(B_\delta\), and at \(V_n=0\) the interval is \(\{\bar X_n\}\).
\end{hbmethod}

\begin{hbexample}{def:p15-w02b-p15-059-magnitude-example}{A negligible departure that is significant, and a larger one that is not}
Take \(m=0\), \(\alpha=0.95\) and a margin \(\delta=0.2\) in the units of
the measurement.

\emph{A large study.} With \(n=100000\), \(\bar X_n=0.021\) and
\(S_n=1.003\): \(e_n\approx0.0031718\), \(T_0\approx6.621\),
\(t_{99999,0.975}\approx1.95999\) and \(\pi_0\approx3.6\times10^{-11}\).
The interval is \(C_0\approx[0.01478,\,0.02722]\). The result is
significant, and \(C_0\subset(0,0.2)\): the interval shows that the
departure is nonzero and that it is negligible, about one tenth of the
margin.

\emph{A small study.} With \(n=20\), \(\bar X_n=0.35\) and \(S_n=1.0\):
\(e_n\approx0.22361\), \(T_0\approx1.565\), \(t_{19,0.975}\approx2.0930\)
and \(\pi_0\approx0.134\). The interval is
\(C_0\approx[-0.118,\,0.818]\); it contains \(0\) and values beyond the
margin, so the result is not significant and the interval is inconclusive
about magnitude. The small study has the larger estimate and the much
larger p-value: the p-values order the two studies opposite to their
estimates.

\emph{Probability of significance for a negligible departure.} Under
normal sampling with \(\sigma=1\) and \(\mu=0.02\), so \(d=0.02\) and
\(|d|<\delta\), the exact probabilities are as follows (to four
decimals).
\begin{center}
\begin{tabular}{rccc}
\(n\) & \(P\{\phi_0^0=1\}\) & \(P(\mathcal N_n)\) & \(P\{\phi_0^0=1,\ C_0\subset(0,0.2)\}\)\\
\hline
\(100\) & \(0.0545\) & \(0.0523\) & below \(10^{-14}\)\\
\(1000\) & \(0.0968\) & \(0.9999\) & \(0.0920\)\\
\(10000\) & \(0.5159\) & \(1.0000\) & \(0.5159\)\\
\(100000\) & \(1.0000\) & \(1.0000\) & \(1.0000\)\\
\end{tabular}
\end{center}
The first column is the two-sided noncentral Student tail; the other two
integrate the normal law of \(\bar X_n\) against the independent law of
\(V_n\) (\cref{thm:p15-student-orthogonal}); all were computed numerically.
At \(n=100\) the half-width \(te_n\) is near \(0.2\), so the interval can
seldom fit inside \(B_\delta\); at \(n=100000\) the departure is significant
with probability \(0.999994\), and the interval places it inside
\((0,0.2)\) with essentially the same probability, as
\cref{thm:p15-w02b-p15-059-magnitude-limits} requires in the limit.

\emph{Posterior statements.} For the large study with \(\sigma=1\) taken as
known and the prior \(\operatorname{normal}(0,1)\),
\(u_n\approx0.021000\) and \(\sqrt{v_n}\approx0.0031623\). Then
\(\Pi_n(\theta\leq0)\approx1.6\times10^{-11}\),
\(\Pi_n(|\theta|\geq0.2)<10^{-697}\), and the central posterior interval
is \(K_n\approx[0.01480,\,0.02720]\). The posterior is close to certain both
that the departure is positive and that it is negligible.

\emph{The scale of \(\sigma\).} At \(n=20\), a departure \(d=0.5\), larger
than the margin, with \(\sigma=10\) has \(\beta_{20}(0.05)\approx0.0552\);
a departure \(d=0.02\) with \(\sigma=0.01\) has
\(\beta_{20}(2)>0.99999999\) (\cref{thm:p15-w02b-p15-059-scale}).
\end{hbexample}

\begin{hbwarning}{def:p15-w02b-p15-059-magnitude-limits}{What significance, confidence and posterior statements say about magnitude}
\begin{enumerate}
\item The error probability of the significance verdict is fixed only at
\(d=0\) (\cref{thm:p15-w02b-p15-059-verdicts}, part~5). A significant
result does not state that \(|d|\) is large, and a non-significant one does
not state that it is small (\cref{thm:p15-w02b-p15-059-scale,%
thm:p15-w02b-p15-059-magnitude-limits}).
\item The verdicts depend on \(\delta\); the same interval gives different
verdicts for different margins. Nothing in this chapter determines
\(\delta\); it has to come from the purpose of the inference.
\item The finite-sample bounds of \cref{thm:p15-w02b-p15-059-verdicts} are
proved under normal sampling. Under finite-variance sampling only the
limits of \cref{thm:p15-w02b-p15-059-magnitude-limits} are proved, and by
\cref{thm:p15-w02b-p15-059-nonuniform} no sample size makes the coverage of
\(C_0\), on which the verdicts rest, uniform over that class.
\item At \(|d|=\delta\) each verdict has limiting probability
\((1-\alpha)/2\); a departure at the margin is not resolved by any sample
size.
\item The posterior limits of \cref{thm:p15-w02b-p15-059-posterior} are for
a fixed true mean, known \(\sigma\) and a normal prior; at a fixed \(n\) the
posterior values depend on \(\mu_0\) and \(\tau^2\). A posterior
probability of a positive departure is not a measure of its size.
\end{enumerate}
\end{hbwarning}

\section{Worked example and limitations}
\label{sec:p15-w02b-p15-059-example}

\begin{hbexample}{def:p15-w02b-p15-059-example}{Normal cutoffs at small \(\nu\), a skewed law, and a uniform endpoint}
\emph{Normal cutoffs.} Take \(\alpha=0.95\), so
\(c=z_{0.975}\approx1.95996\) and \(\phi_0(c)\approx0.058445\). The exact
coverage of \(\bar X_n\pm ce_n\) under normal sampling, and the
approximation \(0.95-c\phi_0(c)(1+c^2)/(2\nu)\) suggested by
\cref{thm:p15-w02b-p15-059-second-order}, are as follows.
\begin{center}
\begin{tabular}{rccc}
\(\nu\) & exact \(\kappa_\nu(c)\) & first-order approximation & \(t_{\nu,0.975}\)\\
\hline
\(1\) & \(2\arctan(c)/\pi\approx0.6997\) & \(0.6727\) & \(12.706\)\\
\(2\) & \(c/\sqrt{2+c^2}\approx0.8109\) & \(0.8114\) & \(4.3027\)\\
\(9\) & \(\approx0.9184\) & \(0.9192\) & \(2.2622\)\\
\(29\) & \(\approx0.9403\) & \(0.9404\) & \(2.0452\)\\
\end{tabular}
\end{center}
The rows for \(\nu=1,2\) use the closed forms of
\cref{thm:p15-student-worked-cases}; those for \(\nu=9,29\) are numerical
values of \(2F_\nu(c)-1\). Every exact entry lies below \(0.95\), as
\cref{thm:p15-w02b-p15-059-mixture-order} requires. The quantile
correction \(z_p+z_p(1+z_p^2)/(4\nu)\) of
\cref{thm:p15-w02b-p15-059-second-order} gives \(3.1461\) at \(\nu=2\)
against the exact \(4.3027\), and \(2.2235\) at \(\nu=9\) against
\(2.2622\): the correction term describes the approach, not the value at
small \(\nu\).

\emph{A skewed law.} With \(n=10\) Bernoulli observations of success
probability \(p=0.01\), all ten are zero with probability
\(0.99^{10}\approx0.9044\). On that event every interval
\(C_k\) is \(\{0\}\) and misses the mean \(0.01\), so every multiplier
gives coverage at most about \(0.0956\). For this fixed law the coverage
still tends to \(0.95\) as \(n\to\infty\)
(\cref{thm:p15-w02b-p15-059-nonnormal}); lowering \(p\) reproduces the
failure at any fixed \(n\) (\cref{thm:p15-w02b-p15-059-nonuniform}).

\emph{A uniform endpoint.} Take \(n=20\), \(\alpha=0.95\) and observed
\(L_{20}=9.8\), \(\bar X_{20}=5.1\). Then
\((0.05)^{-1/20}\approx1.16159\) and \(I_{20}\approx[9.8,\,11.384]\), of
length about \(1.584\). The normal-approximation interval has
\(\widehat\theta=10.2\) and half-width
\(1.95996\times10.2/\sqrt{60}\approx2.581\), so
\(J_{20}\approx[7.619,\,12.781]\), of length about \(5.162\). Its lower
part \([7.619,9.8)\) consists of values already excluded by the observed
maximum.
\end{hbexample}

\begin{hbwarning}{def:p15-w02b-p15-059-limits}{What the large-sample statements do not give}
\begin{enumerate}
\item \cref{thm:p15-w02b-p15-059-nonnormal} is pointwise in the law. By
\cref{thm:p15-w02b-p15-059-nonuniform} no sample size makes the coverage of
\(C_k\) close to \(\alpha\) for all finite-variance laws at once; a
statement of the form ``\(n\geq30\) suffices'' is not a consequence of
either theorem.
\item The correction term of \cref{thm:p15-w02b-p15-059-second-order} is
proved only under normal sampling. For other finite-variance laws this
chapter proves the first-order limit and asserts no rate.
\item \cref{thm:p15-w02b-p15-059-nonnormal} requires finite positive
variance. Nothing is asserted here for laws without a finite variance.
\item The local power limit is proved for the location-shift array of part~5
only; other triangular arrays need a separate central limit argument, for
instance through \cref{thm:p15-clt-lindeberg-feller}, which is not carried
out here.
\item In a nonregular model the regular normal approximation is not the
relevant large-sample regime: \cref{thm:p15-w02b-p15-059-uniform-regimes}
shows a rate \(1/n\) exact interval beside a rate \(n^{-1/2}\) approximate
interval that includes excluded values.
\end{enumerate}
\end{hbwarning}

\emph{Relationships.} The exact regime is \cref{thm:p15-student-law,%
thm:p15-student-coverage,thm:p15-student-two-sided-optimal}; the boundary
cases \(n=2\), \(V=0\) and \(\nu=1\) are in
\cref{thm:p15-student-worked-cases,def:p15-student-model}. The limits here
use \cref{thm:p15-clt-iid} (Durrett's Theorem~3.4.1),
\cref{thm:p15-clt-perturbation}, \cref{thm:p15-weak-continuous-mapping},
\cref{thm:p15-extended-mean-strong-law} and the density argument of
\cref{thm:p15-scheffe-density-limit}. Neyman's two-observation uniform
systems are \cref{def:p15-neyman-uniform-systems,%
thm:p15-neyman-uniform-coverage}; part~1 of
\cref{thm:p15-w02b-p15-059-uniform-regimes} extends the third system to
\(n\) observations. The large-sample normal limit of a robust location
estimator is \cref{thm:p15-convex-location-normality}, and its finite
contamination failures are in \cref{thm:p15-huber-location-failure-cases}.
The p-value of \cref{thm:p15-w02b-p15-059-pvalue} is an instance of
\cref{def:p15-valid-p-value}; the posterior used in
\cref{thm:p15-w02b-p15-059-posterior} is that of
\cref{thm:p15-bayes-normal-calculation}, and the distinction between
coverage and posterior content on which part~5 rests is
\cref{thm:p15-confidence-credibility-boundary,thm:p15-bayes-coverage}. The
verdict \(\mathcal N_n\) is a pair of one-sided Student tests
(\cref{thm:p15-w02b-p15-059-verdicts}, part~2), so its error bound is the
exact size of those tests from \cref{thm:p15-student-coverage}.

\emph{Sources.} The central limit theorem and self-normalized sums:
\cite[p.~144, Theorem~3.4.1; p.~147, Exercise~3.4.5]{DurrettPTEAuthor2019}.
The remaining statements and their proofs are given in the annex.
